\chapter{Podmnožice in potenčne množice}
\label{cha:podmnozice}

V tem poglavju spoznamo podmnožice in potenčno množico, katere elementi so vse podmnožice dane množice.
%
Podmnožico lahko podamo z lastnostjo njenih elementov ali s karakteristično preslikavo.
%
Pokažemo, da sta ta dva opisa enakovredna.
%
Pperacije na podmnožicah povežemo z logičnimi vezniki iz poglavja~\ref{cha:boolova-algebra} in dokažemo, da potenčna množica tvori Boolovo algebro.

\section{Podmnožice}
\label{sec:podmnozice}

Pravimo, da je množica $S$ \df{podmnožica} množice $T$ in pišemo $S \subseteq T$, ko velja $\all{x \in S} x \in T$. Pravimo tudi, da je $S$ \df{vsebovana} v $T$ in da je $T$ \df{nadmnožica}~$S$.
%
Pravilo ekstenzionalnosti za množice lahko zapišemo s~$\subseteq$:
%
\begin{equation*}
  S = T \liff S \subseteq T \land T \subseteq S.
\end{equation*}

\begin{zgled}
  Naj bo $P$ množica vseh praštevil, torej je $P \subseteq \NN$. Podmnožica $P$
  opredeljuje lastnost \nls{je praštevilo}.
\end{zgled}


\begin{zgled}
  Vsaka množica $S$ ima vsaj kako podmnožico, namreč $\emptyset \subseteq S$ in $S \subseteq S$.
  % 
  Vsebovanost $\emptyset \subseteq S$ pomeni $\all{x \in \emptyset} x \in S$, kar je poseben primer zgleda~\ref{zgled:vse-v-prazni}.
\end{zgled}

Podmnožico tvorimo tako, da podamo lastnost, ki pove, kateri elementi nadmnožice so elementi podmnožice:

\begin{pravilo}[Podmnožica]
  \label{pravilo:podmnozica}
  Za vsako množico $A$ in lastnost~$\phi$ njenih elementov lahko tvorimo
  \df{podmnožico}  $\set{x \in A \such \phi(x)}$, katere elementi so
  natanko tisti elementi množice~$A$, ki zadoščajo pogoju~$\phi$:
  %
  \begin{equation*}
    a \in \set{x \in A \such \phi(x)} \liff a \in A \land \phi(a).
  \end{equation*}
\end{pravilo}

\noindent
Pri tem je $x$ vezana spremenljivka.
%
Nekateri avtorji uporabljajo tudi oznako $\set{x \in A \mathbin{;} \phi(x)}$ in redkeje $\set{x \in A \mathbin{,} \phi(x)}$.

Velja $\set{ x \in A \such \phi(x) } \subseteq A$, prav tako pa
\begin{equation*}
  \set{x \in A \such \phi(x)} \subseteq \set{x \in A \such \psi(x)} \liff
  \all{x \in A} \phi(x) \lthen \psi(x).
\end{equation*}

\begin{zgled}
  Lahko bi nas skrbelo, da bi imeli podmnožico $S \subseteq A$, ki je ne moremo tvoriti s pravilom~\ref{pravilo:podmnozica}, ker ne poznamo nobene primerne lastnosti. A to ni res,
  saj lahko vzamemo $\phi(x) \dfeq (x \in S)$, ker velja $S = \set{x \in A \mid x \in S}$.
\end{zgled}

\begin{vaja}
  Denimo, da sta $S$ in $A$ poljubni množici, se pravi, da ne velja nujno $S \subseteq A$ kot v prejšnjem zgledu.
  Kako bi opisali množico $\set{x \in A \mid x \in S}$?
\end{vaja}

Za podmnožico $S \subseteq T$ definiramo \df{kanonično inkluzijo} ali \df{kanonično vključitev}
%
\begin{align*}
   i_S &: S \to T\\
     x &\mapsto x.
\end{align*}
%
Pozor, to ni identiteta, razen v primeru $S = T$!
Oznaka $i_S$ ni standardna, pravzaprav standardne oznake ni.

Če je $f : T \to U$ in $S \subseteq T$, pravimo kompoziciji $f \circ i_S$ \df{zožitev} preslikave $f$ na $S$ in jo pišemo $\restrict{f}{S}$.


\section{Potenčna množica in karakteristične preslikave}
\label{sec:potencna-mnozica}

Vse podmnožice dane množice zberemo v eno množico.

\begin{pravilo}[Potenčna množica]
  \label{pravilo:potencna-mnozica}
  Za vsako množico $A$ lahko tvorimo \df{potenčno množico} $\pow{A}$, katere elementi so natanko podmnožice množice~$A$:
  %
  \begin{equation*}
    S \in \pow{A} \liff S \subseteq A.
  \end{equation*}
\end{pravilo}

Na primer, $\pow{\emptyset} = \set{\emptyset}$ in
%
\begin{equation*}
  \pow{\set{a,b,c}} = \set{\emptyset, \set{a}, \set{b}, \set{c}, \set{a,b}, \set{a,c}, \set{b,c}, \set{a,b,c}}.
\end{equation*}


\df{Karakteristična preslikava} na množici $A$ je preslikava z domeno~$A$ in kodomeno~$\two$, pri čemer je $\two = \set{\bot, \top}$ množica resničnostnih vrednosti.
%
Eksponentna množica $\two^A$ je torej množica vseh karakterističnih preslikav na~$A$.

\begin{opomba}
  Karakteristične preslikave se uporabljajo tudi v analizi, kjer jih
  običajno razumemo kot preslikave $A \to \set{0,1}$. Ker sta množici $\two = \set{\bot,\top}$ in $\set{0,1}$ izomorfni, to ni bistvena razlika.
\end{opomba}

Karakteristična preslikava opredeljuje neko lastnost elementov~$A$: tisti elementi, ki imajo opredeljeno lastnost, se
slikajo v $\top$, ostali pa v $\bot$.

\begin{zgled}
  Preslikava $p : \NN \to \two$, definirana s predpisom
  %
  \begin{equation*}
    p(n) \defeq
    \begin{cases}
      \top & \text{če je $n$ praštevilo}, \\
      \bot & \text{če $n$ ni praštevilo},
    \end{cases}
  \end{equation*}
  %
  je karakteristična preslikava lastnosti \nls{je praštevilo}. Lahko bi jo zapisali tudi takole:
  %
  \begin{equation*}
    p(n) \defeq (n > 1 \land \all{k, m} (n = k \cdot m \lthen (k = 1 \lor m = 1))).
  \end{equation*}
\end{zgled}


\subsection{Izomorfizem \texorpdfstring{$\pow{A} \cong \two^A$}{P(A) ≅ 2^A}}

Videli smo, da lahko neko lastnost elementov množice $A$ predstavimo bodisi s
podmnožico bodisi s karakteristično preslikavo. To nam da idejo, da med
podmnožicami $A$ in karakterističnimi preslikavami na $A$ obstaja neka zveza.

\begin{izrek}
  \label{izrek:potencna-karakteristicne}
  Za vsako množico $A$ velja $\pow{A} \cong \two^A$.
\end{izrek}

\begin{proof}
  Definirajmo preslikavi
  %
  \begin{align*}
    \chi &: \pow{A} \to \two^A &
    \xi &: \two^A \to \pow{A} \\
    \chi_S(x) &\defeq
      \begin{cases}
        \top & \text{če $x \in S$,} \\
        \bot & \text{če $x \not\in S$,}
      \end{cases}
    &
    \xi_f &\defeq \set{x \in A \such f(x) = \top}.
  \end{align*}
  %
  Ta predpisa bi lahko krajše zapisali tudi takole:
  %
  \begin{align*}
  \chi_S(x) &\defeq (x \in S), &
  \xi_f &\defeq \set{x \in A \such f(x) }.
  \end{align*}
  %
  Preslikavi $\chi_S$ pravimo \df{karakteristična preslikava podmnožice $S$}.
  %
  Trdimo, da sta $\chi$ in $\xi$ inverza:
  %
  \begin{enumerate}
  \item 
    Dokažimo $\chi \circ \xi = \id[\two^A]$. Uporabimo pravilo ekstenzionalnosti za preslikave.
    Naj bo $f \in \two^A$. Dokažimo, da je $\chi_{\xi_f} = f$.
    Uporabimo pravilo ekstenzionalnosti za preslikave. Naj bo $x \in A$:
    %
    \begin{equation*}
      \chi_{\xi_f}(x) = (x \in \xi_f) = f(x).
    \end{equation*}

  \item
    Dokažimo $\xi \circ \chi = \id[\pow{A}]$. Uporabimo pravilo ekstenzionalnosti za preslikave. Naj bo $S \in \pow{A}$. Dokažimo, da je $\xi_{\chi_S} = S$:
    %
    \begin{equation*}
      \xi_{\chi_S} = \set{x \in A \such \chi_S(x)} = \set{x \in A \such x \in S} = S.
    \end{equation*}
    %
    Zadnja enakost velja, ker je $S \subseteq A$.
  \end{enumerate}
\end{proof}

\section{Operacije na podmnožicah}
\label{sec:operacije-na-podmnozicah}

Izomorfizem iz izreka~\ref{izrek:potencna-karakteristicne} prenese logične veznike na podmnožice.
%
Naj bo $A$ množica in $S, T \subseteq A$.
%
Definiramo \df{presek}, \df{unijo} in \df{razliko} podmnožic $S$ in~$T$:
%
\begin{align*}
  S \cap T &\defeq \set{x \in A \such x \in S \land x \in T}, \\
  S \cup T &\defeq \set{x \in A \such x \in S \lor x \in T}, \\
  S \setminus T &\defeq \set{x \in A \such x \in S \land \lnot (x \in T)}.
\end{align*}
%
Razliki $A \setminus S$ pravimo \df{komplement} podmnožice~$S$ v~$A$.
%
Unija podmnožic $S$ in $T$ je enaka uniji iz pravila~\ref{pravilo:unija}, saj imata enake elemente.

Karakteristične preslikave teh podmnožic izrazimo z vezniki.
%
Za vsak $x \in A$ velja
%
\begin{align*}
  \chi_{S \cap T}(x) &= \chi_S(x) \land \chi_T(x), &
  \chi_{S \cup T}(x) &= \chi_S(x) \lor \chi_T(x), &
  \chi_{A \setminus S}(x) &= \lnot \chi_S(x).
\end{align*}
%
Poleg tega je $\chi_{\emptyset}(x) = \bot$ in $\chi_A(x) = \top$.
%
Ker je $\chi$ izomorfizem, sta podmnožici enaki natanko tedaj, ko sta enaki njuni karakteristični preslikavi.
%
Zato vsako pravilo Boolove algebre iz razdelka~\ref{sec:boolova-algebra} da pravilo za podmnožice, če v njem zamenjamo $\land$ s~$\cap$, $\lor$ s~$\cup$, $\lnot$ s komplementom, $\bot$ z~$\emptyset$ in $\top$ z~$A$.
%
Na primer, iz distributivnosti $p \land (q \lor r) = (p \land q) \lor (p \land r)$ dobimo
%
\begin{equation*}
  S \cap (T \cup U) = (S \cap T) \cup (S \cap U)
\end{equation*}
%
za vse $S, T, U \subseteq A$.
%
Podmnožice množice $A$ torej tvorijo Boolovo algebro za operacije presek $\cap$, unija $\cup$ in komplement.
%
Nevtralni element za unijo je $\emptyset$ in nevtralni element za presek je $A$.

Definirajmo tudi operacijo \df{simetrična razlika $\oplus$}, ki podmnožicama $S, T \subseteq A$ priredi podmnožico
%
\begin{equation*}
  S \oplus T \defeq (S \setminus T) \cup (T \setminus S) = (S \cup T) \setminus (S \cap T).
\end{equation*}
%
Simetrični razliki ustreza ekskluzivna disjunkcija: za vsak $x \in A$ velja
%
\begin{equation*}
  \chi_{S \oplus T}(x) = \chi_S(x) \veebar \chi_T(x).
\end{equation*}
%
Z resničnostnimi tabelami preverimo, da je ekskluzivna disjunkcija asociativna in komutativna, da je $\bot$ zanjo nevtralni element in da velja $p \veebar p = \bot$.
%
Zato je operacija $\oplus$ asociativna in komutativna, $\emptyset$ je zanjo nevtralni element in velja $S \oplus S = \emptyset$.
%
Potenčna množica $\pow{A}$ je torej za operacijo $\oplus$ komutativna grupa, v kateri je vsak element sam sebi inverz.
